\documentclass[11pt]{article}
\usepackage[margin=1in]{geometry}
\usepackage{amsmath,amssymb,amsthm}
\usepackage{microtype,booktabs}
\usepackage{hyperref}
\hypersetup{colorlinks=true,linkcolor=blue,citecolor=blue,urlcolor=blue,
 pdftitle={Stieltjes moment classification of Apery polynomials},
 pdfauthor={Brandon Li and Ethan Yeroushalmi}}
\setlength{\emergencystretch}{2em}
\newtheorem{theorem}{Theorem}
\newtheorem{lemma}[theorem]{Lemma}
\newtheorem{proposition}[theorem]{Proposition}
\newcommand{\E}{\mathbb E}
\newcommand{\R}{\mathbb R}
\newcommand{\C}{\mathbb C}
\title{Stieltjes moment classification of Ap\'ery polynomials}
\author{Brandon Li \and Ethan Yeroushalmi\\[0.35em]
 \normalsize Independent researchers, Great Neck, New York, USA\\[0.2em]
 \normalsize Corresponding author: Brandon Li,
 \href{mailto:brandon.li.gn@gmail.com}{brandon.li.gn@gmail.com}}
\date{}
\begin{document}
\maketitle
\begin{abstract}
We prove that the Ap\'ery polynomials
$A_n(x)=\sum_{k=0}^n\binom nk^2\binom{n+k}k^2x^k$ form a Stieltjes
moment sequence in $n$, for real $x\ge0$, exactly when $x=0$ or $x\ge1$.
This settles Sokal's conjecture for real parameters; the case $x=1$ is
Edgar's theorem on the Ap\'ery numbers. Both directions start from a
single representation of $A_n(x)$ as the Legendre moments of an explicit
probability law on two angles. For $0<x<1$, a Stieltjes representation
would force the part of this law below $1$ to have a real-analytic
density. The law instead has a support gap when $x<1/4$, a saddle
singularity when $1/4<x<1$, and a degenerate endpoint at $x=1/4$, so no
such representation exists. For $x>1$, the same part of the law has a
density given by a complete elliptic integral. We continue it to the plane
cut along $[1,\infty)$, show that its jump across the cut is positive,
and recover it by Cauchy's formula, using a uniform bound on large
circles. Heine's integral for the Legendre functions $Q_n$ then turns the
remaining Legendre moments into the power moments of a positive measure
on $[0,1]$. We also give exact rational certificates showing that, for
each $0<x\le2/3$, one of four explicit Hankel minors is negative.
\end{abstract}
\noindent\textbf{Keywords:} Ap\'ery polynomials; Stieltjes moment problem;
Legendre polynomials; elliptic integrals; Hankel matrices.

\smallskip\noindent\textbf{Mathematics Subject Classification (2020):}
Primary 44A60; Secondary 33C45, 33E05, 15B05.

\section{Introduction}
A real sequence $(a_n)_{n\ge0}$ is a \emph{Stieltjes moment sequence} if
$a_n=\int_0^\infty t^n\,d\nu(t)$ for a positive Borel measure $\nu$.
Define
\begin{equation}\label{eq:apery}
 A_n(x)=\sum_{k=0}^n\binom nk^2\binom{n+k}k^2x^k.
\end{equation}
Wang and Zhu posed moment questions for generalized Ap\'ery polynomials,
including the family \eqref{eq:apery} \cite[Section 5]{WangZhu}.
Sokal recorded the case $x\ge1$ as an open problem \cite{Sokal2018} and
later conjectured that $(A_n(x))_{n\ge0}$ is a Stieltjes moment sequence
for every $x\ge1$ and is not one for $0<x<1$ \cite[slide 29]{Sokal}.
Edgar proved the case $x=1$ \cite{Edgar}. Our main result is the full
classification for real parameters.
\begin{theorem}\label{thm:classification}
For real $x\ge0$, $(A_n(x))_{n\ge0}$ is a Stieltjes moment sequence
if and only if $x=0$ or $x\ge1$.
\end{theorem}
Sokal's stronger conjecture, that the Hankel matrix of
$(A_n(1+x))_{n\ge0}$ is coefficientwise totally positive as a matrix of
polynomials in $x$, remains open.

\subsection{Outline of the proof}
Throughout, $P_n$ is the Legendre polynomial and $r=\sqrt{x}$. Both
ranges of $x$ are handled through the identity
$A_n(x)=\int P_n(y)\,d\mu_r(y)$ of Section~\ref{sec:average}, where
$\mu_r$ is an explicit probability measure with compact support. The part
of $\mu_r$ on $[1,\infty)$ already gives positive power moments, by
Laplace's integral for $P_n$. Everything therefore depends on the part
on $[-1,1)$.

For $0<r<1$, the Poisson kernel for Legendre series shows that a
Stieltjes representation would force this part to have a real-analytic
density (Section~\ref{sec:regularity}). The geometry of $\mu_r$ rules
this out: its support has a gap when $r<1/2$, it has a saddle
singularity when $r>1/2$, and it has a quartic endpoint when $r=1/2$
(Section~\ref{sec:negative}). For $r>1$, the density is a complete
elliptic integral. We continue it to the plane cut along $[1,\infty)$,
compute its jump across the cut, which is positive, and recover the
density by Cauchy's formula. The classical $Q_n$ transform then turns the
remaining Legendre moments into power moments of a positive measure on
$[0,1]$ (Section~\ref{sec:positive}). Section~\ref{sec:hankel} gives
exact Hankel certificates for $0<x\le2/3$ and explains why no fixed
finite test can cover the whole interval $(0,1)$.

\subsection{Related work}
Common showed that Legendre series with Stieltjes coefficients continue
to a cut plane, and that positive mixtures of the functions $Q_n$ give
Stieltjes moment sequences \cite[CERN preprint, Theorems 2--3 and 5]{Common};
he notes that Yndurain found the second result independently. We give
the finite-measure and endpoint details needed in our setting. The new
ingredients are the geometric obstruction for $0<x<1$ and the elliptic
boundary calculation for $x>1$.

At $x=1$, our angular representation reduces to a formula attributed to
Lupa\c{s} (2005) and proved by Ivanisvili in 2014 \cite{Lupas}; the
product and generating-function identities behind it are classical.
Matala-aho derived a recurrence for the same polynomials
\cite[Theorem~2]{MatalaAho}, and Neuschel studied their asymptotics and
the distribution of their zeros \cite{Neuschel}; we do not need either
result. Zhu and Sun studied arithmetic properties of the Hankel
determinants of the Ap\'ery numbers \cite{ZhuSun}.

\section{A regularity obstruction}\label{sec:regularity}
The Legendre polynomials are normalized by $P_n(1)=1$.
We first record the consequence of the Poisson kernel that we need.
\begin{lemma}\label{lem:regularity}
Let $\mu$ be a finite positive measure supported on $[-1,H]$, where $H\ge1$.
If $a_n=\int P_n(y)\,d\mu(y)$ is a Stieltjes moment sequence, then
$\mu|_{[-1,1)}$ has a density real analytic on $(-1,1)$ and bounded in a
right neighborhood of $-1$.
\end{lemma}
\begin{proof}
For $y\ge1$, Laplace's integral representation gives
\begin{equation}\label{eq:laplace}
 P_n(y)=\frac1\pi\int_0^\pi
 \bigl(y+\sqrt{y^2-1}\cos\phi\bigr)^n\,d\phi.
\end{equation}
Hence the contribution from $[1,H]$ is the moment sequence of a finite positive
measure $\nu_+$ on $[0,R]$, where we may take
$R=\max\{2,H+\sqrt{H^2-1}\}$.
Since $|P_n(y)|\le1$ on $[-1,1]$ and $|P_n(y)|\le R^n$ on $[1,H]$,
we have $|a_n|\le\mu([-1,H])R^n$.
Any positive Stieltjes measure $\nu$ representing $(a_n)$ must therefore be
supported on $[0,R]$: positive mass beyond $R+\delta$, for any $\delta>0$,
would contradict this exponential bound.

The finite signed measure $\sigma=\nu-\nu_+$ has moments
\begin{equation}\label{eq:bn}
 b_n=\int_{[-1,1)}P_n(y)\,d\mu(y),\qquad |b_n|\le\mu([-1,1)).
\end{equation}
These bounds force $\sigma$ to be supported on $[0,1]$, even though $\sigma$
need not be positive. Indeed, its Cauchy transform agrees, for $|z|>R$, with
$\sum_{n\ge0}b_n z^{-n-1}$. The series is holomorphic on $|z|>1$, so by
the identity theorem the two functions agree on
$\{\operatorname{Im}z>0,\ |z|>1\}$, and likewise in the lower
half-plane. The series is real on $(1,\infty)$.
The Stieltjes inversion formula, which holds for finite signed measures by
dominated convergence against $|\sigma|$, therefore gives $\sigma(I)=0$ for
every interval $I\subset(1,\infty)$ whose endpoints are not atoms of $\sigma$.
Such intervals exhaust $(1,\infty)$, so $\sigma$ vanishes there.

For $0\le s<1$ put
\begin{equation}\label{eq:kernel}
 K(s,y)=\frac{1-s^2}{2(1-2sy+s^2)^{3/2}},\qquad -1\le y\le1.
\end{equation}
The generating function for $P_n$, followed by differentiation and
orthogonality, gives
\[
 K(s,y)=\frac12\sum_{n\ge0}(2n+1)s^nP_n(y),\qquad
 \int_{-1}^1P_n(y)K(s,y)\,dy=s^n.
\]
In particular, $K(s,y)\,dy$ is a probability measure. For $s=1$ use
$\delta_1$ instead. Integrating these measures against $\sigma$ produces a
signed measure $\eta$ of total variation at most $|\sigma|([0,1])$.
Fubini's theorem applies because $|P_n|\le1$ on $[-1,1]$, and shows
that the Legendre moments of $\eta$ are the $b_n$.
Since the Legendre polynomials span all polynomials, uniqueness of finite
signed measures on this compact interval implies
$\eta=\mu|_{[-1,1)}$.

Since $\mu|_{[-1,1)}$ has no mass at $1$, the atom
$\sigma(\{1\})\delta_1$ vanishes, and $\mu|_{[-1,1)}$ has density
$\int_{[0,1)}K(s,y)\,d\sigma(s)$.
For each real $y_0<1$, $1-2sy_0+s^2$ has a positive minimum over
$0\le s\le1$. Thus the kernels are holomorphic and uniformly bounded in a
small complex neighborhood of $y_0$. Integration against the finite signed
measure preserves holomorphy. This proves analyticity on $(-1,1)$ and,
by taking $y_0=-1$, boundedness near the left endpoint.
\end{proof}

\section{The angular representation}\label{sec:average}
The classical identities used in this section follow from the Legendre
addition theorem and integral representations; see \cite{DLMF}.
The finite expansion
\[
 P_n(1+2z)=\sum_{k=0}^n\binom nk\binom{n+k}k z^k
\]
and Parseval's identity imply, with $r=\sqrt{x}$,
\begin{equation}\label{eq:parseval}
 A_n(x)=\E_t\left|P_n(1+2r e^{it})\right|^2,
\end{equation}
where $t$ is uniform on $[0,2\pi]$.
The Legendre product formula is
\begin{equation}\label{eq:product}
 P_n(a)P_n(b)=\frac1\pi\int_0^\pi
 P_n\left(ab+\sqrt{a^2-1}\sqrt{b^2-1}\cos\theta\right)d\theta.
\end{equation}
Although initially stated for real arguments, this is a polynomial identity
in $a,b$: after integration, only even powers of the product of square roots
remain. It therefore extends to complex arguments and is independent of the
choices of square roots. For $b=\overline a$, choose conjugate square roots.
Equations \eqref{eq:parseval}--\eqref{eq:product} yield
\begin{equation}\label{eq:average}
 A_n(x)=\E P_n(Y_r),
\end{equation}
where $t,\theta$ are independent and uniform on $[0,2\pi]$, and
\begin{equation}\label{eq:Y}
 Y_r=1+4r^2+4r\cos t
       +4r\sqrt{1+2r\cos t+r^2}\cos\theta.
\end{equation}
(Taking $\theta$ uniform on $[0,2\pi]$ rather than $[0,\pi]$ does not
change the distribution of $\cos\theta$.)
Writing $a=1+2r e^{it}$, we have
$Y_r=|a|^2+|a^2-1|\cos\theta$. Since $|a^2-1|\le|a|^2+1$,
$|a|\le1+2r$ and $|a^2-1|=4r|1+re^{it}|\le4r(1+r)$, we have
\begin{equation}\label{eq:range}
 -1\le Y_r\le1+8r^2+8r.
\end{equation}
For $0<r<1$, the square root in \eqref{eq:Y} is strictly positive and smooth.

When $r=1$, substituting $t=2\alpha$ in \eqref{eq:Y}
gives the same distribution as
$1+8\cos^2\alpha+8\cos\alpha\sin\beta$ with independent uniform angles:
the absolute value of $\cos\alpha$ is absorbed by symmetry in the second
angle. This is the formula of Lupa\c{s} \cite{Lupas}; the proof does not
use it.

\section{Failure on \texorpdfstring{$0<x<1$}{0<x<1}}\label{sec:negative}
\begin{theorem}\label{thm:main}
For every fixed real $x$ with $0<x<1$, the sequence
$(A_n(x))_{n\ge0}$ is not a Stieltjes moment sequence.
\end{theorem}
\begin{proof}[Proof of Theorem~\ref{thm:main}]
Fix $0<r<1$, and suppose that $(A_n(x))_{n\ge0}$ is a Stieltjes moment
sequence.
Lemma~\ref{lem:regularity} applies to the distribution $\mu_r$ of $Y_r$.
Set
\[
 h=1+8r^2-8r.
\]
At $(t,\theta)=(\pi,\pi)$, $Y_r=h<1$, so by continuity the law has
positive mass below $1$. We distinguish three cases.

\medskip\noindent\textbf{Case 1: $0<r<1/2$.}
For $c=\cos t$, the lower endpoint as $\theta$ varies is
\[
 f(c)=1+4r^2+4rc-4r\sqrt{1+2rc+r^2}.
\]
Since
\[
 f'(c)=4r\left(1-\frac{r}{\sqrt{1+2rc+r^2}}\right)>0,
 \qquad -1\le c\le1,
\]
where $\sqrt{1+2rc+r^2}\ge1-r>r$, the global lower endpoint is
$f(-1)=h\in(-1,1)$. Thus $\mu_r$ vanishes on $(-1,h)$ but has positive
mass in $(h,1)$. The analytic density furnished by Lemma~\ref{lem:regularity}
would vanish identically on the connected interval $(-1,1)$, a contradiction.

\medskip\noindent\textbf{Case 2: $1/2<r<1$.}
Let $u=t-\pi$ and $v=\theta-\pi$. Expanding about this interior point of
the torus gives
\begin{equation}\label{eq:hessian}
 Y_r=h+\frac{2r(1-2r)}{1-r}u^2+2r(1-r)v^2
        +O\bigl((|u|+|v|)^4\bigr).
\end{equation}
The two quadratic coefficients are nonzero and have opposite signs.
The Morse lemma supplies local smooth coordinates, with a bounded positive
Jacobian, in which $Y_r-h=U^2-V^2$. After a further invertible linear change,
this is the product of the coordinates. In a fixed small square the region
where the absolute value of that product is at most $\epsilon$ has area
at least $c\epsilon\log(1/\epsilon)$ for all sufficiently small
$\epsilon>0$. Since $t$ and $\theta$ are uniformly distributed, it follows
that
\[
 \mu_r([h-\epsilon,h+\epsilon])
       \ge c'\epsilon\log(1/\epsilon).
\]
Since $h\in(-1,1)$, this contradicts the $O(\epsilon)$ bound required
by a density bounded near $h$.

\medskip\noindent\textbf{Case 3: $r=1/2$.}
Now $h=-1$ and direct expansion gives
\begin{equation}\label{eq:quartic}
 Y_{1/2}=-1+\tfrac12(u^4+v^2)
                 +O(u^6+u^2v^2+v^4).
\end{equation}
Together with \eqref{eq:range}, this implies
$0\le Y_{1/2}+1\le C(u^4+v^2)$ locally.
For a sufficiently small fixed $c>0$, the rectangle
\[
 |u|\le c\epsilon^{1/4},\qquad |v|\le c\epsilon^{1/2}
\]
is mapped into $[-1,-1+\epsilon]$. It has probability at least
$c'\epsilon^{3/4}$. This contradicts the $O(\epsilon)$ mass bound
from the density bounded near $-1$ in Lemma~\ref{lem:regularity}.
The three cases cover $0<x=r^2<1$.
\end{proof}


\section{A positive elliptic kernel for \texorpdfstring{$x>1$}{x>1}}
\begin{theorem}\label{thm:positive}
For every real $x>1$, the sequence $(A_n(x))_{n\ge0}$ has a positive
compactly supported Stieltjes representing measure.
\end{theorem}
Fix $r=\sqrt{x}>1$ and set
\[
 h=1+8r^2-8r,\qquad H=1+8r^2+8r,\qquad c_r=\frac1{2\pi^2r}.
\]
Write $f_r$ for the density of the Legendre averaging law on $(-1,1)$.
We will continue $f_r$ to $\Omega=\C\setminus[1,\infty)$, express it as
the Cauchy transform of a positive density, and convert its Legendre
moments into power moments.

\subsection{The density and its continuation}
Throughout this section $\mathcal K$ uses a \emph{parameter}, not a modulus:
\[
 \mathcal K(w)=\int_0^{\pi/2}(1-w\sin^2\theta)^{-1/2}\,d\theta
             =R_F(0,1-w,1).
\]
Here $R_F$ is Carlson's symmetric integral, continued from
\[
 R_F(a,b,c)=\frac12\int_0^\infty
              \frac{dt}{\sqrt{(t+a)(t+b)(t+c)}}.
\]
Each factor uses its principal square root on the initial domain.
These conventions agree with \cite[Sections 19.16 and 19.25]{DLMF}
after replacing its squared modulus by our parameter.

For $-1<y<1$, put
\begin{equation}\label{eq:AB}
 d=\sqrt{(y+1)/2},\quad m=\frac{y-1}{4r},\quad
 A=(1-d)^2-m^2,\quad B=(1+d)^2-m^2 .
\end{equation}
Conditional on $c=\cos t$, the second angular variable has an arcsine
law. Its squared radicand factors as
\[
 16r^2(1+2rc+r^2)-(y-1-4r^2-4rc)^2
 =16r^2\bigl[d^2-(c-m)^2\bigr].
\]
The conditional law has no atoms, since its amplitude is at least
$4r(r-1)>0$. Integrating against the arcsine density of $c$ gives
\[
 f_r(y)=\frac1{4\pi^2r}\int_{m-d}^{m+d}
 \frac{dc}{\sqrt{(1-c^2)(c-m+d)(m+d-c)}} .
\]
Indeed, the roots of the conditional quadratic radicand are $m-d,m+d$,
both strictly inside $(-1,1)$ when $r>1$ and $-1<y<1$.
For example,
$1+m-d=(1-d)(1-(d+1)/(2r))>0$ and
$1-m-d=(1-d)(1+(d+1)/(2r))>0$.
The four-root substitution
\[
 t^2=\frac{(e_3-e_1)(c-e_2)}{(e_3-e_2)(c-e_1)},
 \qquad (e_1,e_2,e_3,e_4)=(-1,m-d,m+d,1),
\]
reduces this integral to the standard complete elliptic integral. Since
$(e_3-e_1)(e_4-e_2)=B$ and
$(e_3-e_2)(e_4-e_1)=4d$, it yields
\begin{equation}\label{eq:ellipticf}
 f_r(y)=c_r R_F(0,A,B)
       =c_r B^{-1/2}\mathcal K(4d/B).
\end{equation}

\begin{lemma}\label{lem:directanalytic}
The germ in \eqref{eq:ellipticf} extends to a single-valued holomorphic
function on $\Omega$.
\end{lemma}
\begin{proof}
On the two-sheeted surface $d^2=(z+1)/2$, the finite singularities of the
composition in \eqref{eq:ellipticf} can occur only above
$-1,1,h,H$. Indeed, $\mathcal K(p)=\pi\,{}_2F_1(1/2,1/2;1;p)/2$
continues along every path avoiding $p=0,1,\infty$, by the classical
hypergeometric equation \cite[15.10.1]{DLMF}. Here $p=4d/B$, and the
preimages are $d=0$, $A=0$, or $B=0$. The factorizations
\begin{equation}\label{eq:ABfactor}
 A=(d-1)^2\left(1-\frac{(d+1)^2}{4r^2}\right),\qquad
 B=(d+1)^2\left(1-\frac{(d-1)^2}{4r^2}\right)
\end{equation}
give precisely those projected points. Critical points of the parameter
map create no singularity in a composition, and the prefactor square root
continues along paths avoiding its zeros. Thus the selected germ continues
along every path in $\Omega\setminus\{-1\}$.

The apparent interior puncture is removable for this germ. The change
$c=m+d\cos\phi$ in the density integral gives, near $d=0$,
\[
 f_r=\frac1{4\pi^2r}\int_0^\pi
 \left[1-\left(\frac{d^2-1}{2r}+d\cos\phi\right)^2\right]^{-1/2}d\phi.
\]
At $d=0$ the expression inside the inner square has modulus
$1/(2r)<1$. For sufficiently small complex $d$ the integrand has a
single holomorphic square-root branch uniformly in $\phi$. Its integral
is holomorphic in $d$, and $\phi\mapsto\pi-\phi$ shows it is even.
It is therefore holomorphic in $z=2d^2-1$ and agrees with the initial germ
for real $d>0$. The generator loop about $-1$ has trivial monodromy.
Since $\Omega$ is simply connected and the fundamental group of
$\Omega\setminus\{-1\}$ is generated by that loop, the monodromy theorem
gives single-valuedness. The local integral then fills the puncture.
\end{proof}

\subsection{The boundary values}
The parameter versions of
the classical connection identities are
\begin{align}
 \mathcal K(z)&=(1-z)^{-1/2}\mathcal K\bigl(z/(z-1)\bigr),
                 && z<0,\label{eq:imaginaryK}\\
 \mathcal K(z)&=z^{-1/2}\bigl[
     \mathcal K(1/z)+i\mathcal K(1-1/z)\bigr],
                 && z\text{ above }(1,\infty).\label{eq:reciprocalK}
\end{align}
They hold as identities of the adjacent holomorphic germs; the lower-bank
version of \eqref{eq:reciprocalK} is its conjugate. See
\cite[19.7.2--19.7.3]{DLMF}, including its explicit upper/lower sign rule.
The convergent logarithmic expansion \cite[19.12.1]{DLMF} gives
\begin{equation}\label{eq:Klog}
 \mathcal K(1-\epsilon)
   =-\frac{\mathcal K(\epsilon)}{\pi}\log\epsilon+G(\epsilon),
\end{equation}
where $G$ is holomorphic near zero.

For real $y>1$, retain the positive square root in $d$.
The factorizations \eqref{eq:ABfactor} show that $A$ has a double zero
at $1$ and a simple zero at $h$, while
$B$ has a simple zero at $H$.
\begin{lemma}\label{lem:upperbank}
The upper boundary value of $f_r$ is given by the following holomorphic
germs adjacent to the indicated real intervals:
\begin{align}
 f_r(z)&=\frac{c_r}{\sqrt B}
   \bigl[\mathcal K(p)+2i\mathcal K(1-p)\bigr],
 &p=\frac{4d}{B},\quad&1<y<h,\label{eq:bankone}\\
 f_r(z)&=\frac{c_r}{\sqrt{4d}}
   \bigl[\mathcal K(q)+3i\mathcal K(1-q)\bigr],
 &q=\frac{B}{4d},\quad&h<y<H,\label{eq:banktwo}\\
 f_r(z)&=\frac{c_r}{\sqrt{-A}}
   \bigl[-2\mathcal K(t)+3i\mathcal K(1-t)\bigr],
 &t=\frac BA,\quad&H<y<\infty.\label{eq:bankthree}
\end{align}
The parameters $p,q,t$ respectively lie in $(0,1)$ on their displayed
real intervals, and all displayed real prefactors use positive roots.
\end{lemma}
\begin{proof}
Near $1$,
\[
 A=\frac{1-r^{-2}}{16}(z-1)^2+O((z-1)^3),\qquad B=4+O(z-1).
\]
Going around $1$ through the upper half-plane makes $A/B$ wind clockwise
once about zero. Equation~\eqref{eq:Klog} adds exactly
$2i\mathcal K(A/B)$, proving \eqref{eq:bankone} near $1$.
On compact subintervals of $(1,h)$, both $p$ and $1-p$ stay in $(0,1)$.
Overlapping neighborhoods therefore propagate the same germ without an
additional period.

At $h$,
\[
 A'(h)=-\frac{(r-1)^2}{r(2r-1)}<0,\qquad
 p'(h)=-A'(h)/B(h)>0 .
\]
An upper semicircle about $h$ thus sends $p$ through the upper side of $1$.
Apply \eqref{eq:reciprocalK} to $\mathcal K(p)$ and
\eqref{eq:imaginaryK} to $\mathcal K(1-p)$:
\[
 \mathcal K(1-p)=p^{-1/2}\mathcal K(1-1/p).
\]
Substitution in \eqref{eq:bankone} gives \eqref{eq:banktwo}, including
the coefficient $3$. The square-root factors agree on these germs because
$Bp=4d$ and they agree before crossing. Again, compact subintervals of
$(h,H)$ keep $q,1-q$ strictly between $0$ and $1$.

At $H$,
\[
 B'(H)=-\frac{(r+1)^2}{r(2r+1)}<0,\qquad q'(H)=B'(H)/(4d(H))<0.
\]
Hence $q$ crosses zero through the lower half-plane and $1-q$ crosses $1$
through the upper half-plane. Apply \eqref{eq:imaginaryK} to $\mathcal K(q)$
and \eqref{eq:reciprocalK} to $\mathcal K(1-q)$. With
\[
 \frac{q}{q-1}=t,\qquad \frac1{1-q}=1-t,\qquad
 (4d)(1-q)=-A,
\]
the real coefficient becomes $1-3=-2$, giving \eqref{eq:bankthree}.
The parameters $t,1-t$ remain in $(0,1)$ to the right of $H$.
Overlapping neighborhoods propagate this germ to every finite point of
that bank. These are valid continuations of the selected germ in
$\Omega$, where Lemma~\ref{lem:directanalytic} supplies single-valuedness.
\end{proof}

Define $\rho_r(y)=\operatorname{Im}f_r(y+i0)/\pi$.
Lemma~\ref{lem:upperbank} gives the explicitly positive density
\begin{equation}\label{eq:rho}
 \rho_r(y)=\frac1{2\pi^3r}
 \begin{cases}
 2R_F(0,4d,B),&1<y<h,\\
 3R_F(0,4d,B),&h<y<H,\\
 3R_F(0,-A,-B),&H<y<\infty.
 \end{cases}
\end{equation}
Thus $\rho_r$ has a finite right limit at $1$, a finite jump at $h$,
and at most a logarithmic singularity at $H$. At infinity,
$-A,-B\sim y^2/(16r^2)$, so homogeneity gives $\rho_r(y)=O(1/y)$.

\subsection{The Cauchy representation}
\begin{lemma}\label{lem:dispersion}
For every $z\in\Omega$,
\begin{equation}\label{eq:dispersion}
 f_r(z)=\int_1^\infty\frac{\rho_r(a)}{a-z}\,da .
\end{equation}
\end{lemma}
\begin{proof}
The integral converges by the local logarithmic bounds and
$\rho_r(a)=O(1/a)$. To apply Cauchy's theorem, we need a bound that is
uniform around a large circle, including angles close to the cut.

Take $R_0=1000r^2$ and the connected upper exterior domain
$D=\{z:\operatorname{Im}z>0,\ |z|>R_0\}$. Choose
$d=\sqrt{(z+1)/2}$ by continuation from the positive bank. The expansion
\[
 A=-\frac{z^2}{16r^2}
   +\left(\frac12+\frac1{8r^2}\right)z
   +\frac32-\frac1{16r^2}-2\sqrt{(z+1)/2}
\]
gives $|A+z^2/(16r^2)|\le5R/4$, $R=|z|>R_0$, at every angle. Thus
\[
 \frac{R^2}{32r^2}\le |A|\le\frac{3R^2}{32r^2},\qquad
 \frac{\sqrt R}{2}\le |d|\le\sqrt R .
\]
Put $w=1-t=-4d/A$ and $A=-z^2(1+e)/(16r^2)$. Then
\[
 \frac{64r^2}{3R^{3/2}}\le |w|
       \le\frac{128r^2}{R^{3/2}}<\frac12,\qquad |e|\le\frac1{50}.
\]
Use $0<\arg z<\pi$ and the small analytic logarithms of $1+1/z$ and
$1+e$. The logarithm of $w$ continued from the positive bank is
\begin{equation}\label{eq:continuedlog}
 \operatorname{Log}w=\log\frac{64r^2}{\sqrt2}
 -\frac32\operatorname{Log}z+\frac12\log(1+1/z)-\log(1+e).
\end{equation}
In particular,
\[
 |\operatorname{Im}\operatorname{Log}w|
 \le\frac{3\pi}{2}+\tfrac12\arcsin(1/1000)+\arcsin(1/50)<5,
 \qquad |\operatorname{Log}w|\le\tfrac32\log R+5.
\]
Here $R^{-3/2}<|w|<1$ controls the real part. This logarithm is
\emph{continued}, not replaced by the principal logarithm: its imaginary
part can approach $-3\pi/2$ near the negative real axis.
The square root of $-A$ on $D$ is chosen positive at the positive bank;
its modulus is $|A|^{1/2}$.

In \eqref{eq:Klog}, the convergent series \cite[19.12.1--19.12.3]{DLMF} is
\[
 G(w)=\sum_{j\ge0}\left(\frac{(1/2)_j}{j!}\right)^2
       [\psi(j+1)-\psi(j+1/2)]w^j.
\]
Its coefficients satisfy $((1/2)_j/j!)^2\le1$ and
$0<\psi(j+1)-\psi(j+1/2)\le2\log2$. Hence for $|w|\le1/2$,
\[
 |\mathcal K(w)|\le\pi/\sqrt2,\qquad |G(w)|\le4\log2.
\]
Equation~\eqref{eq:bankthree} becomes
\begin{equation}\label{eq:fullcircle}
 f_r(z)=\frac{c_r}{\sqrt{-A}}
 \left[\frac2\pi\mathcal K(w)\operatorname{Log}w
       -2G(w)+3i\mathcal K(w)\right].
\end{equation}
All terms on the right are holomorphic on $D$ with the specified branches.
The identity theorem extends this identity from its positive-bank
neighborhood to all $D$. It gives the explicit full-semicircle bound
\[
 |f_r(z)|\le\frac{\sqrt{32}}{2\pi^2R}
 \left[\sqrt2\left(\tfrac32\log R+5\right)
             +8\log2+\frac{3\pi}{\sqrt2}\right].
\]
Schwarz reflection gives the lower semicircle, and holomorphy supplies the
negative-axis limit. Thus the estimate is uniform on the entire slit
circle, and the large-circle Cauchy integral vanishes.

At $1$, $A/B$ has a double zero and \eqref{eq:Klog} gives
$f_r(z)=O(1+|\log|z-1||)$. At $h$, $1-p$ has a simple zero in
\eqref{eq:bankone}; at $H$, $q$ has a simple zero in
\eqref{eq:banktwo}. The same logarithmic expansion therefore gives
$f_r(z)=O(1+|\log|z-a||)$ near each $a\in\{1,h,H\}$ on either
side of the cut. An arc of radius $\epsilon$ contributes
$O(\epsilon(1+|\log\epsilon|))$, which tends to zero.
Cauchy's theorem
on the slit plane now gives \eqref{eq:dispersion}, with
$(f_+-f_-)/(2\pi i)=\rho_r$ and denominator $a-z$.
\end{proof}

\subsection{The positive remainder measure}
Common's positive Froissart--Gribov theorem
\cite[CERN preprint, Section 4, Theorem 5]{Common} converts positive
$Q_n$ mixtures into Stieltjes moments. We give the integral calculation
for the present density, whose spectral support reaches $1$ and whose
total spectral mass need not be finite.

Use the standard normalization
\[
 Q_n(a)=\frac12\int_{-1}^1\frac{P_n(y)}{a-y}\,dy,\qquad a>1.
\]
Heine's integral is
\[
 Q_n(a)=\int_0^\infty
       \bigl(a+\sqrt{a^2-1}\cosh u\bigr)^{-n-1}\,du.
\]
Set $t=(a+\sqrt{a^2-1}\cosh u)^{-1}$. Since
$du=-dt/(t\sqrt{1-2at+t^2})$, this gives
\begin{equation}\label{eq:Qpositive}
 Q_n(a)=\int_0^{a-\sqrt{a^2-1}}
       \frac{t^n}{\sqrt{1-2at+t^2}}\,dt .
\end{equation}
In \cite[14.12.13--14.12.14]{DLMF}, Olver's normalized function is
$Q_n/n!$; multiplying both formulas by $n!$ yields precisely this
normalization.

\begin{proof}[Proof of Theorem~\ref{thm:positive}]
First take $n=0$. All kernels in \eqref{eq:dispersion} and
\eqref{eq:Qpositive} are nonnegative on their real integration domains,
so Tonelli's theorem gives
\[
 \int_{-1}^1 f_r(y)\,dy
 =2\int_1^\infty\rho_r(a)Q_0(a)\,da
 =\int_0^1 w_r(t)\,dt<\infty,
\]
where $w_r$ is defined below. The first integral is the probability
$\mu_r([-1,1))$; the endpoints have zero mass by the conditional
arcsine representation. For general $n$, put
$b_n=\int_{-1}^1P_n(y)f_r(y)\,dy$. The bounds $|P_n(y)|\le1$ and
$0\le t^n\le1$ reduce absolute integrability to the case $n=0$.
Fubini's theorem now gives
\begin{align}
 b_n&=2\int_1^\infty \rho_r(a)Q_n(a)\,da
      =\int_0^1 t^n w_r(t)\,dt,\label{eq:positiveremainder}\\
 w_r(t)&=2\int_1^{(t+t^{-1})/2}
       \frac{\rho_r(a)}{\sqrt{1-2at+t^2}}\,da\ge0,
            \qquad 0<t<1.\label{eq:wmeasure}
\end{align}
The domain in \eqref{eq:wmeasure} follows from
$t\le a-\sqrt{a^2-1}$ if and only if
$1\le a\le(t+t^{-1})/2$. Thus $w_r(t)\,dt$ is a finite positive
measure on $[0,1]$ with moments $b_n$.

The contribution from $Y_r\ge1$ already has the positive Laplace-mixture
measure $\nu_+$, supported on
$[H-\sqrt{H^2-1},H+\sqrt{H^2-1}]$.
Adding the positive measure in \eqref{eq:wmeasure} proves that
$\nu_++w_r(t)\,dt$ is a positive compact Stieltjes measure with moments
$A_n(x)$.
\end{proof}

\begin{proof}[Proof of Theorem~\ref{thm:classification}]
Theorem~\ref{thm:main} excludes $0<x<1$ and
Theorem~\ref{thm:positive} covers $x>1$. Edgar's prior theorem
\cite{Edgar} covers $x=1$. At $x=0$, $A_n(0)=1$ is represented by
the point mass at $1$.
\end{proof}


\section{Hankel certificates and fixed-size tests}\label{sec:hankel}
We now give exact sign certificates on an interval of parameters. Define
\[
 H_{d,s}(x)=\det\bigl(A_{i+j+s}(x)\bigr)_{0\le i,j<d},\qquad s\in\{0,1\}.
\]
\begin{proposition}\label{prop:finite}
For every $0<x\le2/3$, at least one of
$H_{4,0}(x),H_{8,0}(x),H_{9,1}(x),H_{10,0}(x)$ is strictly negative.
\end{proposition}
\begin{proof}
Table~\ref{tab:certificates} gives exact sign certificates. For each row,
$Q(x)=H_{d,s}(x)/x^v$ is a polynomial of degree $N$.

\begin{table}[ht]
\centering
\caption{Intervals and polynomial degrees for the negative-minor certificates.}
\label{tab:certificates}
\begin{tabular}{ccccc}
\toprule
$(d,s)$ & $a$ & $b$ & $v$ & $N$\\[2pt]
\midrule
$(4,0)$ & $0$ & $9/100$ & $3$ & $9$\\
$(8,0)$ & $1/20$ & $101/500$ & $14$ & $42$\\
$(9,1)$ & $201/1000$ & $47/100$ & $18$ & $63$\\
$(10,0)$ & $41/100$ & $2/3$ & $23$ & $67$
\\\bottomrule
\end{tabular}
\end{table}

Writing $Q(a+(b-a)t)=\sum_{j=0}^N c_jt^j$, its Bernstein coefficients are
\[
 \beta_k=\sum_{j=0}^k c_j\frac{\binom{k}{j}}{\binom{N}{j}},
 \qquad 0\le k\le N.
\]
In each of the four rows, every $\beta_k$ is a strictly negative rational
number. The supplementary material records all coefficients; its
verifier recomputes the determinants over $\mathbb Q[x]$, verifies their
valuations and degrees, checks each rational sign, and reconstructs
\[
 Q(a+(b-a)t)=\sum_{k=0}^N\beta_k\binom Nk t^k(1-t)^{N-k}.
\]
The Bernstein basis is nonnegative and sums to one on $[0,1]$, proving
$Q<0$ throughout $[a,b]$. Hence $H_{d,s}<0$ on that interval wherever
$x>0$. The four intervals overlap and cover $(0,2/3]$.
\end{proof}

\subsection{Finite tests near \texorpdfstring{$x=1$}{x=1}}
No fixed finite set of Hankel minors detects every parameter in
$(0,1)$. Indeed, Edgar's theorem \cite[Theorem 1]{Edgar} provides a nonzero
positive integrable moment density at $x=1$. For any fixed $d$ and
$s\in\{0,1\}$, and any nonzero real vector $v$,
\[
 \sum_{i,j=0}^{d-1}v_i v_j A_{i+j+s}(1)
 =\int_0^\infty t^s\left(\sum_{i=0}^{d-1}v_i t^i\right)^2d\nu(t)>0.
\]
A nonzero polynomial has finitely many zeros, whereas the absolutely
continuous measure cannot be supported there. Thus each fixed Hankel
matrix is positive definite at $x=1$, and by continuity any finite
collection of them, or of their principal submatrices, stays positive
definite in a neighborhood of $1$. At each $x\in(0,1)$,
Theorem~\ref{thm:main} gives a negative principal minor of the ordinary
or shifted Hankel matrix, but it does not say which minor or bound its
size.

\subsection{Shifted Hankel matrices}
\begin{proposition}\label{prop:shifts}
For every fixed $x>0$ and positive integer $d$, the matrix
$(A_{N+i+j}(x))_{0\le i,j<d}$ is positive definite for all sufficiently
large integers $N$.
\end{proposition}
\begin{proof}
Use \eqref{eq:average} and the Laplace representation \eqref{eq:laplace}
to write $A_n(x)=a_n+b_n$, where $a_n=\int t^n\,d\nu_+(t)$ is the
positive moment contribution from $Y\ge1$, and $|b_n|\le1$ is the
contribution from $-1\le Y<1$. Put $r=\sqrt{x}$,
$H=1+8r^2+8r>1$ and $R=H+\sqrt{H^2-1}>1$.
The law of $Y$ has positive mass arbitrarily close to its maximum $H$,
attained at $t=\theta=0$ in \eqref{eq:Y}. In the Laplace representation,
varying its independent cosine angle gives $\nu_+$ infinite support
arbitrarily close to $R$. Choose $1<c<R$ so that $\nu_+|_{[c,R]}$
has infinite support. Its degree-$(d-1)$ Gram matrix
\[
 G=\left(\int_{[c,R]}t^{i+j}\,d\nu_+(t)\right)_{0\le i,j<d}
\]
is strictly positive definite, since a nonzero polynomial cannot vanish
on that support. Write $\lambda=\lambda_{\min}(G)>0$.
The matrix $(a_{N+i+j})$ dominates $c^NG$ in positive-semidefinite order,
whereas $(b_{N+i+j})$ has operator norm at most $d$. Hence
\[
 \lambda_{\min}\bigl((A_{N+i+j}(x))_{0\le i,j<d}\bigr)
 \ge c^N\lambda-d>0
\]
for all sufficiently large $N$.
\end{proof}
Thus fixed-size tests at increasingly large shifts become positive on
both sides of the parameter threshold. The cutoff may depend on $x$ and $d$.

\section*{Supplementary information}
The supplementary material contains Python scripts and the exact rational
certificates for Proposition~\ref{prop:finite}, checks of the algebraic
identities used in the proofs, and high-precision numerical checks of the
elliptic continuation and the dispersion formula.

\section*{Statements and Declarations}
\paragraph{Funding.} No funding was received for this work.

\paragraph{Competing interests.} The authors have no relevant financial or
non-financial interests to disclose.

\paragraph{Author contributions.} Brandon Li carried out the research and
prepared the manuscript. Ethan Yeroushalmi checked the proofs and
computations and edited the text. Both authors read and approved the final
manuscript.

\paragraph{Use of AI tools.} The authors used Claude Opus~5.5 (Anthropic),
through Claude Code version 2.1.317, to explore proof strategies, carry out
and check derivations and symbolic computations, write the verification
code, search the literature, and draft and edit the text. The authors take
full responsibility for the content, including the proofs, computations and
references.

\paragraph{Data availability.} The verification scripts and certificate
data are available at \url{https://doi.org/10.5281/zenodo.24022053} and as
supplementary material. No other data were used.

\begin{thebibliography}{10}
\providecommand{\doi}[1]{\url{https://doi.org/#1}}

\bibitem{Common}
Common, A.K.: Properties of Legendre expansions related to series of
Stieltjes and applications to $\pi$--$\pi$ scattering. Nuovo Cimento A
\textbf{63}, 863--891 (1969). \doi{10.1007/BF02760744}. Preprint
version: CERN TH.1034 (1969).
\url{https://cds.cern.ch/record/348575/files/CM-P00058257.pdf}

\bibitem{Lupas}
Edgar, G.A.: A conjectured formula for Ap\'ery numbers. MathOverflow,
question 178790 (2014), with an answer by P.~Ivanisvili.
\url{https://mathoverflow.net/questions/178790}. Accessed 7 November 2026

\bibitem{Edgar}
Edgar, G.A.: The Ap\'ery numbers as a Stieltjes moment sequence.
Preprint at \url{https://arxiv.org/abs/2005.10733} (2020)

\bibitem{MatalaAho}
Matala-aho, T.: On recurrences for some hypergeometric type polynomials.
Ast\'erisque \textbf{198--200}, 237--244 (1991)

\bibitem{Neuschel}
Neuschel, T.: Ap\'ery polynomials and the multivariate saddle point
method. Constr. Approx. \textbf{40}, 487--507 (2014).
\doi{10.1007/s00365-014-9245-3}

\bibitem{DLMF}
Olver, F.W.J., Olde Daalhuis, A.B., Lozier, D.W., Schneider, B.I.,
Boisvert, R.F., Clark, C.W., Miller, B.R., Saunders, B.V., Cohl, H.S.,
McClain, M.A. (eds.): NIST Digital Library of Mathematical Functions.
Release 1.2.8 of 2026-09-15. \url{https://dlmf.nist.gov/}

\bibitem{Sokal2018}
Sokal, A.D.: The Euler and Springer numbers as moment sequences. Expo.
Math. \textbf{38}, 1--26 (2020). \doi{10.1016/j.exmath.2018.08.001}

\bibitem{Sokal}
Sokal, A.D.: Stieltjes moment sequences and coefficientwise Hankel-total
positivity in enumerative combinatorics. Lecture at Guttmann80, 1 July
2025. \url{https://clisby.net/guttmann2025/talks/Alan_Sokal.pdf}.
Accessed 7 November 2026

\bibitem{WangZhu}
Wang, Y., Zhu, B.-X.: Log-convex and Stieltjes moment sequences. Adv.
Appl. Math. \textbf{81}, 115--127 (2016). \doi{10.1016/j.aam.2016.06.008}

\bibitem{ZhuSun}
Zhu, B.-X., Sun, Z.-W.: Hankel-type determinants for some combinatorial
sequences. Int. J. Number Theory \textbf{14}, 1265--1277 (2018).
\doi{10.1142/S1793042118500793}

\end{thebibliography}
\end{document}
